% content/c04-quantentheorie/en.tex — Chapter 4: Quantum theory as a probabilistic theory
% (English edition; the German de.tex leads, same constructs, IDs and displays)
%
% Scope: the material of the sixth lecture (latex/lectures/de/vl06.tex, "Postulates of quantum
% mechanics"), 2026-11-04, decision D-C4 in docs/plan/lec56/BRIEF.md: core in finite dimension
% with the qubit examples, deeper blocks only for the hazards H4.1 (normal and singular states)
% and H4.3 (self-adjoint Hamiltonian).  Gleason and Busch, Naimark, instruments and complete
% positivity come later.
% Where every line of vl06 went: evidence/migrations/maps/vl06-map.json.
% The postulates P0 to P5 are TJO's wording from docs/plan/BRIEF.md:50-55 (ledger c04.cl.p0 to
% c04.cl.p5, author-stipulated); the English branch of the displays is that wording verbatim
% (without the marker "(*)" of P1 and with the course's symbols: F for E, \BdOp, \Id).  The German
% branch is a translation by the chapter author, not TJO's wording (GAP-TJO in de.tex).
% Derivations: evidence/derivations/c04/DERIVATION.md, numbers checked by
% evidence/derivations/c04/check_c04.py.
% Voice: sentences of vl06.tex are kept where they are right; remarks from the videos
% (transcripts/clean/v07.md, v08.md up to 44:00) are commented with their time stamp.
% Hats (CONV-HAT): effect operators, projectors, Hamiltonian, unitary operators and observables
% carry the hat; states rho, the identity \Id, ket-bras and the Pauli matrices do not.  The letter
% of the effect operator is open (collision COL-F, BRIEF.md writes E, vl06 F); \EffOp sets F for
% now.

\chapter{Quantum theory as a probabilistic theory}\label{c04}

% Voice: vl06.tex:3-5; v07 [03:49]-[05:36] ("classically, probability is always an expression of
% missing information").
The numbers that quantum mechanics delivers as results are probabilities or, somewhat more
generally, expectation values. That the theory delivers only probabilities at all is the real
break with classical physics. Classically a probability is always an expression of missing
information; in principle everything there can be as precise as you like. Quantum mechanically
it is fundamental, and there are measurements that cannot be carried out jointly with full
precision.

% Voice: vl06.tex:7-13; v07 [05:36]-[10:00].  Shortened: confidence intervals, the lab next door.
One compares a probability with experiment by repeating the experiment sufficiently often and
forming the relative frequency of hits; an expectation value, by averaging the measured values
(Chapter~\ref{c01}). Every run carried out by the same procedure, within given tolerances, counts
as a repetition of ``the same'' experiment; it need not be the same apparatus. In every experiment
one distinguishes two parts, the \emph{preparation} and the \emph{measurement}
(Chapter~\ref{c00}, \defref{D1}). The theory must therefore contain at least three things:
quantities describing the preparations, quantities describing the measurements, and a formula
that assigns a probability to every combination.

Exactly these three ingredients were introduced in Chapter~\ref{c02} for an arbitrary
probabilistic theory: a state space, the effects as affine functions on it, and the value of an
effect in a state as a probability. Quantum theory is such a theory. This chapter says which one:
its states and effects are built from a Hilbert space, with the tools of Chapter~\ref{c03}.

\section{The postulates}\label{c04.s1}

% Voice: v07 [10:00]-[11:56] ("the statement that we do not understand quantum mechanics is
% wrong"; "every book chooses a different collection").  vl06.tex:17-19 (greybox).
\begin{greybox}[On the status of the postulates]
  You often hear that we do not understand quantum mechanics. That misses the point: we have
  mastered the formalism, and for over fifty years we have had postulates that give us the best
  models of the universe. What we cannot do is \emph{justify} them any further; they are given to
  us by nature. Every textbook chooses a slightly different collection of postulates. At bottom
  they are all equivalent, but showing the equivalence takes work.
\end{greybox}

We set up the postulates step by step, each more specialised than the one before. You know the first
from Chapter~\ref{c00}:
% P0 repeats c00.eq.p0 word for word; P1 to P5 are worded as in script v0.2 (lec06.tex, TJO's
% text; decision D-WORTLAUT in docs/plan/lec56/BRIEF.md, 2026-11-06).
\begin{equation}\label{c04.eq.p0}\claim{c04.cl.p0}{author-stipulated}
  \begin{aligned}
    &\text{P0.\ \tpcifenglish{Physical theories are probabilistic theories:}{Physikalische Theorien sind probabilistische Theorien:}} \\
    &\text{\tpcifenglish{predictions in the form of probabilities.}{Vorhersagen in Form von Wahrscheinlichkeiten.}}
  \end{aligned}
\end{equation}
The next postulate fixes what the theory is built from.
\begin{equation}\label{c04.eq.p1}\claim{c04.cl.p1}{author-stipulated}
  \begin{aligned}
    &\text{P1.\ \tpcifenglish{Every kind of quantum mechanical system is assigned a Hilbert space}{Jeder Sorte quantenmechanischer Systeme wird ein Hilbertraum}}\ \Hil\text{\tpcifenglish{.}{ zugeordnet.}}
  \end{aligned}
\end{equation}

% Voice: vl06.tex:24-34; v07 [15:58] ("one postulate is still missing": several particles),
% [19:49] (completeness), [23:01] (wave function).  BRIEF.md:57 ((*): "Many-body QM is an
% extension of this axiom, and introduces the tensor operation, direct sum.").  The example
% L^2(R^3) stands as a pointer; the core stays finite-dimensional (decision D4).
For a \emph{qubit} $\Hil = \mathbb{C}^2$, for a \emph{qudit} $\Hil = \mathbb{C}^d$. For a particle
in space it is the space $\Ltwo{\mathbb{R}^3}$ of square-integrable functions; which Hilbert space
belongs to a system is determined by experiment, not by the theory, and Chapter~5 takes up this
question for a particle on a line. P1 speaks of \emph{one} system. The Hilbert space of several
systems is built from those of the parts, with a further recipe, the tensor product; that comes in
Chapter~12. In the core of this chapter $\Hil$ is finite-dimensional, as in Chapter~\ref{c03}.

\subsection{States}\label{c04.s1.1}

\begin{equation}\label{c04.eq.p2}\claim{c04.cl.p2}{author-stipulated}
  \begin{aligned}
    &\text{P2.\ \tpcifenglish{A preparation procedure is characterised by a state, or density matrix.}{Ein Präparierverfahren wird charakterisiert durch einen Zustand, bzw. eine Dichtematrix.}} \\
    &\text{\tpcifenglish{This is a positive operator}{Das ist ein positiver Operator}}\ \rho\ \text{\tpcifenglish{on}{auf}}\ \Hil\text{\tpcifenglish{, with}{, mit}}\ \tr(\rho) = 1 . \\
    &\text{\tpcifenglish{For pure states in particular}{Speziell für reine Zustände ist}}\ \rho\ \text{\tpcifenglish{is a one-dimensional projector}{ein eindimensionaler Projektor}}\ \rho = \ketbra{\psi}{\psi} , \\
    &\text{\tpcifenglish{hence given by a vector}{also gegeben durch einen Vektor}}\ \ket{\psi} \in \Hil .
  \end{aligned}
\end{equation}
% Voice: vl06.tex:36-41 (postulate 2, "state, or density matrix").  D22 repeats the notion of the
% postulate because the registry needs it for \defref.
\tpcdef{D22}

Chapter~\ref{c02} introduced states more abstractly, as elements of a state space (\defref{D3})
on which the effects are affine functions. P2 says which set that is for quantum theory: the
density operators. One obtains the same set by introducing a state as a rule that assigns to
every quantity its expectation value, that is, as a positive linear functional with norm~1 on the
set $\BdOp{\Hil}$ of operators on $\Hil$. A \emph{linear functional} $\AlgSt$ assigns a number
$\AlgSt(\op{A})$ to every operator $\op{A}$, linearly; it is called \emph{positive} if
$\AlgSt(\op{A}) \geq 0$ for every positive operator $\op{A}$, and for a positive functional
``norm~1'' is the same as $\AlgSt(\Id) = 1$. In finite dimension every such rule is given by a
density operator that one evaluates with the trace:

% Not in vl06: the equivalence with the brief's version (BRIEF.md:52, positive normalised
% functionals).  The chapter author's derivation (DERIVATION.md, c04.cl.funktionale).
\begin{proposition}[States are density operators]\label{c04.prop.funktionale}\claim[proof-included]{c04.cl.funktionale}{derived}
  Let $\Hil$ be finite-dimensional. For every positive linear functional $\AlgSt$ on the operators
  with $\AlgSt(\Id) = 1$ there is exactly one density operator $\rho$ with
  \begin{equation}\label{c04.eq.funktional}
    \AlgSt(\op{A}) = \tr(\rho\, \op{A}) \qquad \text{\tpcifenglish{for every operator}{für jeden Operator}}\ \op{A} .
  \end{equation}
  Conversely, for every density operator $\rho$ the rule $\op{A} \mapsto \tr(\rho\,\op{A})$ is
  such a functional.
\end{proposition}
\begin{proof}
  Let $\{\ket{x}\}$ be an orthonormal basis. We set $\melem{y}{\rho}{x} = \AlgSt(\ketbra{x}{y})$;
  this fixes an operator $\rho$. By Proposition~\ref{c03.prop.ketbra},
  $\op{A} = \sum_{x,y} \melem{x}{\op{A}}{y} \ketbra{x}{y}$, so by linearity
  $\AlgSt(\op{A}) = \sum_{x,y} \melem{x}{\op{A}}{y}\melem{y}{\rho}{x} = \tr(\rho\,\op{A})$. For every
  vector, $\ketbra{\psi}{\psi}$ is positive, so with Proposition~\ref{c03.prop.spur}
  $\melem{\psi}{\rho}{\psi} = \tr(\rho\,\ketbra{\psi}{\psi}) = \AlgSt(\ketbra{\psi}{\psi}) \geq 0$:
  $\rho$ is positive. Further, $\tr(\rho) = \AlgSt(\Id) = 1$. Two operators with the same trace
  against every $\ketbra{x}{y}$ have the same matrix elements; so $\rho$ is unique. Conversely,
  $\op{A} \mapsto \tr(\rho\,\op{A})$ is linear, $\tr(\rho\,\Id) = 1$, and for positive $\op{A}$
  with spectral decomposition $\op{A} = \sum_x a_x \ketbra{\phi_x}{\phi_x}$, $a_x \geq 0$,
  $\tr(\rho\,\op{A}) = \sum_x a_x \melem{\phi_x}{\rho}{\phi_x} \geq 0$.
\end{proof}

In finite dimension this says everything about the states. In infinitely many dimensions
$\tr(\rho) = 1$ demands more, namely that $\rho$ belongs to the trace class, and there are
positive normalised functionals that come from no density operator\zoomref{c04.v1}.

\begin{vertiefung}[id=c04.v1,
    titel={States in infinitely many dimensions},
    wozu={In the core the positive normalised functionals are exactly the density operators
      (Proposition~\ref{c04.prop.funktionale}). P2 demands $\tr(\rho) = 1$; in infinitely many
      dimensions that presupposes the trace class, and there are states that come from no density
      operator. This block identifies the states represented by density operators and gives a
      state that is not represented by one.},
    voraussetzung={Section~\ref{c04.s1.1}; blocks~\ref{c03.v1} and~\ref{c03.v5} (bounded
      operators). The core does not need this block.},
    zurueck={The density operators are the normal states. In finite dimension these are all
      states; in infinitely many dimensions there are others, and the course does not use them.}]
  % Hazard H4.1 (docs/plan/research/curriculum/REPORT.md); GAP-TJO G-0003: whether the course
  % introduces the states as all positive normalised functionals on B(H), as the normal ones or as
  % both is TJO's decision.  This block puts both side by side.
  In infinitely many dimensions $\BdOp{\Hil}$ is the set of bounded operators
  (block~\ref{c03.v5}), and the \emph{norm} $\norm{\op{T}}$ of such an operator is the smallest
  number $C$ with $\norm{\op{T}\psi} \leq C\norm{\psi}$ for all $\ket{\psi}$. A linear functional
  $\AlgSt$ has the norm $\norm{\AlgSt} = \sup\{\abs{\AlgSt(\op{T})} : \norm{\op{T}} \leq 1\}$.
  For a positive functional $\norm{\AlgSt} = \AlgSt(\Id)$ \srcref{bratteli-robinson-1}{p.~49};
  ``positive with norm~1'' therefore means, as in the core, $\AlgSt(\Id) = 1$. A
  \emph{density operator} is now a positive operator $\rho$ for which the series
  $\sum_x \melem{\phi_x}{\rho}{\phi_x}$ over an orthonormal basis converges and has the value $1$;
  these are the positive trace-class operators with trace $1$. As in finite dimension, every
  density operator gives a state $\op{A} \mapsto \tr(\rho\,\op{A})$. The following theorem of
  Bratteli and Robinson identifies the states that arise in this way; here a state is called
  \emph{normal} if it is compatible with monotonically increasing families of positive operators
  that are bounded above, $\AlgSt(\sup_\alpha \op{A}_\alpha) = \sup_\alpha \AlgSt(\op{A}_\alpha)$,
  where $\sup_\alpha \op{A}_\alpha$ is the least upper bound.

  \begin{theorem}[Normal states]\label{c04.thm.normal}\claim[proof-omitted]{c04.cl.normal}{source-quoted}
    A state $\AlgSt$ on $\BdOp{\Hil}$ is normal if and only if there is a density operator $\rho$
    with $\AlgSt(\op{A}) = \tr(\rho\,\op{A})$ for all $\op{A}$.
  \end{theorem}

  Bratteli and Robinson prove the theorem for every von Neumann algebra, in particular for
  $\BdOp{\Hil}$ \srcref{bratteli-robinson-1}{p.~76}. That in infinitely many dimensions there are
  states that are not normal is shown by the following construction.

  \begin{example}[A state without a density operator]\label{c04.ex.singulaer}\claim{c04.cl.singulaer}{derived}
    Let $\Hil = \elltwo$. On the operators of the form $c\Id + \op{F}$, where $\op{F}$ has only
    finitely many linearly independent image vectors, we set $\AlgSt_0(c\Id + \op{F}) = c$. This is
    well defined, because $\Id$ has infinitely many linearly independent image vectors. For a
    vector $\ket{v}$ of norm $1$ with $\op{F}\ket{v} = 0$, which exists because $\op{F}$ has an
    infinite-dimensional kernel, $\norm{(c\Id + \op{F})v} = \abs{c}$; so
    $\abs{\AlgSt_0(c\Id + \op{F})} \leq \norm{c\Id + \op{F}}$, and $\AlgSt_0$ has norm $1$. By the
    Hahn-Banach theorem, $\AlgSt_0$ can be extended to a linear functional $\AlgSt$ on all of
    $\BdOp{\Hil}$ with the same norm \srcref{reed-simon-1}{p.~77}. Because
    $\AlgSt(\Id) = 1 = \norm{\AlgSt}$, $\AlgSt$ is positive, hence a state. It gives every
    one-dimensional projector the value $\AlgSt(\ketbra{\psi}{\psi}) = 0$. If it had a density
    operator $\rho$, then $\melem{\phi_x}{\rho}{\phi_x} = 0$ for every vector of an orthonormal
    basis, so $\tr(\rho) = 0$, a contradiction.
  \end{example}

  Such a state gives the probability zero to every detector that ``sees'' only finitely many
  dimensions. It is not a vector state on this Hilbert space and not a finite or countable mixture
  of such. The argument for its existence is not constructive: the Hahn-Banach extension is not
  given explicitly. For systems with finitely many particles the course uses only the normal
  states, that is, the density operators; that is the trace class which $\tr(\rho) = 1$ in P2
  presupposes.
\end{vertiefung}

% Voice: vl06.tex:43-45; v07 [24:47]-[32:10] ("wave functions are not the states"; "there are no
% closed systems in physics").  Hazard H3.2.
A warning is in order here: the Hilbert space does \emph{not} directly describe the states. Many
books say the states are the wave functions; that is not quite right. The Hilbert space is the
space \emph{in which} we define states; the states themselves are operators. The wave functions
describe only the pure states. A system that interacts with its environment is in general not in
a pure state, and completely isolated systems do not exist in physics: the electromagnetic field
always couples in. Mixed states are therefore the normal case, not a concession to mathematical
generality.

% P2 (pure states = one-dimensional projectors); D4 (Chapter 2): pure = extreme point.
% vl06.tex:47-52 (one-dimensional projector, normalisation through the trace); v07 [32:10] (rank).
The density operators form a state space in the sense of \defref{D3}: mixtures
$\lambda\rho_1 + (1-\lambda)\rho_2$ with $0 \leq \lambda \leq 1$ are again density operators.
P2 calls the one-dimensional projectors the pure states. That these are exactly the states that
are not proper mixtures (\defref{D4}) is the content of the following proposition.

\begin{proposition}[Pure states]\label{c04.prop.rein}\claim[proof-included]{c04.cl.rein}{derived}
  Let $\Hil$ be finite-dimensional. A density operator is pure if and only if it has the form
  $\rho = \ketbra{\psi}{\psi}$ with $\norm{\psi} = 1$. Every density operator is a mixture of pure
  states: with its spectral decomposition,
  \begin{equation}\label{c04.eq.mischung}
    \rho = \sum_x p_x \ketbra{\phi_x}{\phi_x} ,
    \qquad p_x \geq 0 , \quad \sum_x p_x = 1 .
  \end{equation}
\end{proposition}
\begin{proof}
  By Theorem~\ref{c03.thm.spektral}, $\rho$ has a spectral decomposition; the eigenvalues are not
  negative by Equation~\eqref{c03.eq.orte}, and their sum is $\tr(\rho) = 1$. That is
  Equation~\eqref{c04.eq.mischung}. If $\rho$ has two eigenvalues $p_x > 0$, then
  Equation~\eqref{c04.eq.mischung} is a proper mixture of different states, and $\rho$ is not pure.
  So a pure state has the form $\ketbra{\phi}{\phi}$. Conversely, let
  $\ketbra{\psi}{\psi} = \lambda\rho_1 + (1-\lambda)\rho_2$ with $0 < \lambda < 1$. For every
  $\ket{\phi}$ with $\braket{\psi}{\phi} = 0$ we then have
  $0 = \lambda\melem{\phi}{\rho_1}{\phi} + (1-\lambda)\melem{\phi}{\rho_2}{\phi}$, and both terms
  are not negative, hence zero. Writing $\rho_1 = \sum_x q_x \ketbra{\xi_x}{\xi_x}$, this means
  $\sum_x q_x \abs{\braket{\xi_x}{\phi}}^2 = 0$: every eigenvector $\ket{\xi_x}$ with $q_x > 0$
  is orthogonal to all vectors orthogonal to $\ket{\psi}$, and is therefore a multiple of
  $\ket{\psi}$. Hence $\rho_1 = c\ketbra{\psi}{\psi}$, and $\tr(\rho_1) = 1$ gives $c = 1$;
  likewise $\rho_2$. The normalisation: by Equation~\eqref{c03.eq.spur-ketbra},
  $\tr(\ketbra{\psi}{\psi}) = \braket{\psi}{\psi}$, and this is $1$ if and only if
  $\norm{\psi} = 1$.
\end{proof}

A pure state is thus given by a normalised vector $\ket{\psi}$, its \emph{wave function};
$\ket{\psi}$ and $e^{i\alpha}\ket{\psi}$ give the same state. One also calls $\ketbra{\psi}{\psi}$
a one-dimensional projector. The word ``one-dimensional'' refers to no spatial dimension, but to
the dimension of the image: the eigenspace with eigenvalue $1$ is one-dimensional. The usual
phrase ``the system is in the state $\ket{0}$'' is convenient but imprecise: a state characterises
a preparation procedure, not a single system.

% Voice: vl06.tex:54-61; v07 [33:00]-[39:03] ("three real numbers").  The connection to the ball
% of Chapter 2 is a calculation of the chapter author (DERIVATION.md, c04.cl.qubit-zustaende).
\begin{example}[States of a qubit]\label{c04.ex.qubit-zustaende}\claim{c04.cl.qubit-zustaende}{derived}
  For $\Hil = \mathbb{C}^2$, $\rho$ is a positive $2\times 2$ matrix with trace one. Positivity
  requires $\rho_{00} = p \geq 0$, the trace condition forces $\rho_{11} = 1-p$, and because
  $\rho_{11}\geq 0$ must hold as well, $p$ lies between zero and one. Because $\rho$ is Hermitian,
  the off-diagonal entries are complex conjugates of each other:
  \begin{equation}\label{c04.eq.qubit}
    \rho = \begin{pmatrix} p & x \\ x^* & 1-p \end{pmatrix} .
  \end{equation}
  This does not yet exhaust positivity: it requires that both eigenvalues be non-negative, and
  here that means
  \begin{equation}\label{c04.eq.det}
    \det\rho = p(1-p) - \abs{x}^2 \geq 0 ,
    \qquad \text{\tpcifenglish{that is}{das heißt}} \qquad
    \bigl(p - \tfrac12\bigr)^2 + \abs{x}^2 \leq \tfrac14 .
  \end{equation}
  A qubit state thus depends on three real numbers, $p$, $\real x$ and $\imag x$, and they lie in
  a ball of radius $\tfrac12$. With $r_z = 2p - 1$ and $x = \tfrac12(r_x - i r_y)$ this is the
  ball of states $\norm{\vb{r}} \leq 1$ of Proposition~\ref{c02.prop.bloch}. The pure states lie
  on the boundary, $\det\rho = 0$.
\end{example}

\begin{exercise}[The eigenvalues]\label{c04.exc.eigenwerte}\claim{c04.cl.eigenwerte}{derived}
  Compute the two eigenvalues of the matrix~\eqref{c04.eq.qubit} and derive from them the
  condition~\eqref{c04.eq.det}.
\end{exercise}
\begin{solution}
  The eigenvalues are the zeros of
  $\det(\rho - \lambda\Id) = \lambda^2 - \lambda + p(1-p) - \abs{x}^2$, that is
  $\lambda_\pm = \tfrac12 \pm \sqrt{\tfrac14 - p(1-p) + \abs{x}^2}$. Their product is $\det\rho$
  and their sum $1$. Both are non-negative if and only if their product is non-negative, that is
  $p(1-p) - \abs{x}^2 \geq 0$. Because $p(1-p) = \tfrac14 - (p - \tfrac12)^2$, this is the second
  form.
\end{solution}

\subsection{Measurements}\label{c04.s1.2}

\begin{equation}\label{c04.eq.p3}\claim{c04.cl.p3}{author-stipulated}
  \begin{aligned}
    &\text{P3.\ \tpcifenglish{A detector, that is a measuring apparatus with the two possible responses}{Ein Detektor, das heißt ein Messapparat mit den beiden Ansprechmöglichkeiten}} \\
    &\text{\tpcifenglish{``yes'' (hit) and ``no'', is given by an operator}{„ja“ (Treffer) und „nein“, ist gegeben durch einen Operator}}\ \op{\EffOp}\ \text{\tpcifenglish{with}{mit}}\ 0 \leq \op{\EffOp} \leq \Id .
  \end{aligned}
\end{equation}
A detector is thus a yes-no question in the sense of Chapter~\ref{c02}.
% Voice: vl06.tex:66; v07 [52:53]-[57:08] ("operator inequalities are a generalisation"; "the one
% is the identity operator").
The notation $0 \leq \op{\EffOp} \leq \Id$ needs explaining. For operators, $\op{A} \leq \op{B}$
means that $\op{B} - \op{A}$ is positive. So P3 requires \emph{two} conditions,
$\op{\EffOp} \geq 0$ \emph{and} $\Id - \op{\EffOp} \geq 0$, and on the right stands the identity
operator, not the number one.

\tpcdef{D24}

Where exactly these operators come from is said by Chapter~\ref{c02}: there an effect is an
affine function on the state space with values between $0$ and $1$ (\defref{D2}); the detectors
are thus the dual of the state space. The effects of the state space of the density operators
are exactly the effect operators:

% Not in vl06; the chapter author's derivation (DERIVATION.md, c04.cl.effekte).  GAP-TJO G-0011:
% whether the no-restriction hypothesis (D12) is part of P3; the proposition shows only that the
% effects of the state space are the effect operators.
\begin{proposition}[Effects are effect operators]\label{c04.prop.effekte}\claim[proof-included]{c04.cl.effekte}{derived}
  Let $\Hil$ be finite-dimensional.
  \begin{enumerate}
    \item For a Hermitian operator $\op{\EffOp}$, $\tr(\rho\,\op{\EffOp})$ lies between $0$ and $1$
      for every density operator $\rho$ if and only if $0 \leq \op{\EffOp} \leq \Id$.
    \item Every affine function $f$ on the density operators with values between $0$ and $1$ has
      the form $f(\rho) = \tr(\rho\,\op{\EffOp})$ with exactly one effect operator $\op{\EffOp}$.
  \end{enumerate}
\end{proposition}
\begin{proof}
  For~1: with $\rho = \ketbra{\psi}{\psi}$, $\tr(\rho\,\op{\EffOp}) = \melem{\psi}{\op{\EffOp}}{\psi}$
  (Equation~\eqref{c03.eq.zyklisch}); if this lies in $[0, 1]$ for every normalised vector, then
  $\op{\EffOp}$ and $\Id - \op{\EffOp}$ are positive. Conversely, with
  Equation~\eqref{c04.eq.mischung},
  $\tr(\rho\,\op{\EffOp}) = \sum_x p_x \melem{\phi_x}{\op{\EffOp}}{\phi_x}$ is a mixture of numbers
  in $[0, 1]$. For~2: we extend $f$ to the positive operators, by $g(0) = 0$ and
  $g(\op{P}) = \tr(\op{P})\, f\bigl(\op{P}/\tr(\op{P})\bigr)$ for $\op{P} \neq 0$. Because $f$ is
  affine, $g(\op{P} + \op{Q}) = g(\op{P}) + g(\op{Q})$ and $g(c\op{P}) = c\, g(\op{P})$ for
  $c \geq 0$. By its spectral decomposition every Hermitian operator is a difference of two positive
  ones, $\op{A} = \op{P} - \op{Q}$, and $g(\op{A}) = g(\op{P}) - g(\op{Q})$ does not depend on the
  choice: from $\op{P} - \op{Q} = \op{P}' - \op{Q}'$ it follows that
  $\op{P} + \op{Q}' = \op{P}' + \op{Q}$. Every operator is $\op{A} = \op{A}_1 + i\op{A}_2$ with
  Hermitian $\op{A}_1 = \tfrac12(\op{A} + \op{A}\adj)$ and
  $\op{A}_2 = \tfrac{1}{2i}(\op{A} - \op{A}\adj)$, and $\AlgSt(\op{A}) = g(\op{A}_1) + i\,g(\op{A}_2)$
  is linear. As in the proof of Proposition~\ref{c04.prop.funktionale} there is exactly one operator
  $\op{\EffOp}$ with $\AlgSt(\op{A}) = \tr(\op{\EffOp}\,\op{A})$; it is Hermitian because $\AlgSt$
  is real on Hermitian operators. On the density operators $\AlgSt = f$, and by~1,
  $0 \leq \op{\EffOp} \leq \Id$.
\end{proof}

% Voice: vl06.tex:69-80; v07 [57:08]-[67:09] ("the Z measurement can really be carried out";
% "qubits model systems with angular momentum").
\begin{example}[The Z measurement on a qubit]\label{c04.ex.z-messung}\claim{c04.cl.z-messung}{derived}
  The standard measurement on today's quantum computers is the Z measurement, the detector
  \begin{equation}\label{c04.eq.z}
    \op{\EffOp} = \ketbra{0}{0} = \begin{pmatrix} 1 & 0 \\ 0 & 0 \end{pmatrix} .
  \end{equation}
  Both conditions have to be checked. $\op{\EffOp}$ is a projector and hence positive
  (Proposition~\ref{c03.prop.positiv}). Separately, $\Id - \op{\EffOp} = \ketbra{1}{1}$ is a
  projector as well, hence positive. Only now may we call $\op{\EffOp}$ a detector. Physically,
  qubits often describe an angular momentum, for instance a spin~$\tfrac12$; if $\ket{0}$ is the
  state with $z$-component $+\hbar/2$, this detector asks whether the measurement of the
  $z$-component gives this value.
\end{example}

\subsection{The Born rule}\label{c04.s1.3}

\begin{equation}\label{c04.eq.p4}\claim{c04.cl.p4}{author-stipulated}
  \begin{aligned}
    &\text{P4.\ \tpcifenglish{The response frequency of a detector}{Die Ansprechhäufigkeit eines Detektors}}\ \op{\EffOp}\ \text{\tpcifenglish{on a system prepared according to}{an einem gemäß}}\ \rho\ \text{\tpcifenglish{is}{präparierten System ist}}\ \tr(\rho\, \op{\EffOp}) . \\
    &\text{\tpcifenglish{In a pure state}{In einem reinen Zustand}}\ \ket{\psi}\ \text{\tpcifenglish{this expression specialises to}{spezialisiert sich dieser Ausdruck zu}}\ \melem{\psi}{\op{\EffOp}}{\psi} .
  \end{aligned}
\end{equation}

\tpcdef{D26}

% vl06.tex:83 (addition on pure states) and 99-104 (calculation with cyclicity).
The addendum about pure states is not a further assumption, but a calculation: with
$\rho = \ketbra{\psi}{\psi}$ and the cyclicity of the trace, Equation~\eqref{c03.eq.zyklisch},
\begin{equation}\label{c04.eq.born-rein}\claim{c04.cl.born-rein}{derived}
  \tr\bigl(\ketbra{\psi}{\psi}\op{\EffOp}\bigr) = \tr\bigl(\bra{\psi}\op{\EffOp}\ket{\psi}\bigr) = \melem{\psi}{\op{\EffOp}}{\psi} .
\end{equation}
If $\op{\EffOp} = \ketbra{\phi}{\phi}$ is a one-dimensional projector, this is
$\abs{\braket{\phi}{\psi}}^2$; in this form the Born rule appears in most textbooks.
% Voice: v07 [67:09] ("here postulate 2, here postulate 3, here postulate 4").
Make it a habit at the start to say explicitly, in every calculation, which postulate you are
using.

\begin{example}[Detection frequency of the Z measurement]\label{c04.ex.born}\claim{c04.cl.born}{derived}
  With the state~\eqref{c04.eq.qubit} and the detector~\eqref{c04.eq.z},
  \begin{equation}\label{c04.eq.born-z}
    \rho\,\op{\EffOp} = \begin{pmatrix} p & 0 \\ x^* & 0 \end{pmatrix} ,
    \qquad
    \tr(\rho\,\op{\EffOp}) = p .
  \end{equation}
  By P4 we get a hit with probability $p$. This justifies reading the diagonal entries of a
  density matrix as probabilities, namely for the measurement in the basis in which one writes
  the matrix down.
\end{example}

% Voice: vl06.tex:106-118; v08 [00:00]-[20:00] (observable as a family F(S)).  vl06 calls the
% family "observable"; the core says "measurement", as D9 does in Chapter 2 (G-0074 open).
A measuring apparatus often delivers more than a yes or no, namely one of several outcomes $x$
from an outcome set $\Outcomes$. This can be reduced to yes-no questions: for every outcome one
asks whether it has occurred. The corresponding effect operators $\op{\EffOp}_x$ form a
measurement in the sense of \defref{D24}, and the Born rule gives a probability distribution on
$\Outcomes$:
\begin{equation}\label{c04.eq.verteilung}\claim{c04.cl.verteilung}{derived}
  \tr(\rho\,\op{\EffOp}_x) \geq 0 ,
  \qquad
  \sum_{x \in \Outcomes} \tr(\rho\,\op{\EffOp}_x) = \tr\Bigl(\rho \sum_{x} \op{\EffOp}_x\Bigr) = \tr(\rho) = 1 .
\end{equation}
This is Equation~\eqref{c02.eq.messung} of Chapter~\ref{c02}. The question ``does the outcome lie
in the subset $S$?'' has the effect operator $\sum_{x \in S} \op{\EffOp}_x$. For $S = \Outcomes$
this is $\Id$: some outcome occurs for certain; that is why one also includes ``no hit'' as an
outcome in $\Outcomes$.

\begin{example}[A detector with an efficiency]\label{c04.ex.wirkungsgrad}\claim{c04.cl.wirkungsgrad}{derived}
  A detector that reports the state $\ket{0}$ only with probability $q$, $0 < q < 1$, and never
  responds in the state $\ket{1}$, has the effect operator $\op{\EffOp} = q\ketbra{0}{0}$: its
  eigenvalues $q$ and $0$ lie between $0$ and $1$. Together with $\Id - \op{\EffOp}$ it is a
  measurement with two outcomes, and in the state~\eqref{c04.eq.qubit} it responds with
  probability $q p$. It is not a projector, because $\op{\EffOp}^2 = q^2 \ketbra{0}{0} \neq \op{\EffOp}$.
\end{example}

% vl06.tex:118-124; hazard H4.4: "observables are Hermitian operators" is the projection-valued
% special case.  The converse via the spectral decomposition is the chapter author's derivation.
If the outcomes are real numbers, their expectation value (\defref{D55}) is
\begin{equation}\label{c04.eq.erwartung}\claim{c04.cl.erwartung}{derived}
  \sum_{x \in \Outcomes} x\, \tr(\rho\,\op{\EffOp}_x) = \tr(\rho\, \op{A}) ,
  \qquad
  \op{A} = \sum_{x \in \Outcomes} x\, \op{\EffOp}_x ,
\end{equation}
and $\op{A}$ is Hermitian. In the projection-valued case the converse holds too: the measurement
is already determined by $\op{A}$.

\begin{proposition}[Projection-valued measurements and Hermitian operators]\label{c04.prop.pvm}\claim[proof-included]{c04.cl.pvm}{derived}
  Let $\Hil$ be finite-dimensional.
  \begin{enumerate}
    \item If $\{\op{P}_x\}_{x \in \Outcomes}$ is a projection-valued measurement with different
      real outcomes $x$, then $\op{P}_x\op{P}_y = 0$ for $x \neq y$, and $\op{A} = \sum_x x\,\op{P}_x$
      is Hermitian with the eigenvalues $x$ for which $\op{P}_x \neq 0$; $\op{P}_x$ is the
      projector onto the eigenspace for $x$.
    \item Conversely, for every Hermitian operator $\op{A}$ with the different eigenvalues
      $\lambda$, the family of projectors $\op{P}_\lambda$ onto the eigenspaces is a
      projection-valued measurement with $\op{A} = \sum_\lambda \lambda\,\op{P}_\lambda$.
  \end{enumerate}
\end{proposition}
\begin{proof}
  For~1: let $\ket{v} = \op{P}_x\ket{v}$. Then
  $\braket{v}{v} = \sum_y \melem{v}{\op{P}_y}{v} = \braket{v}{v} + \sum_{y \neq x} \norm{\op{P}_y v}^2$,
  because $\melem{v}{\op{P}_y}{v} = \braket{\op{P}_y v}{\op{P}_y v}$ for a projector. So
  $\op{P}_y\ket{v} = 0$ for $y \neq x$, and because every $\op{P}_x\ket{w}$ is such a $\ket{v}$,
  it follows that $\op{P}_y\op{P}_x = 0$. Hence $\op{A}\op{P}_x = x\,\op{P}_x$: every non-zero
  vector in the image of $\op{P}_x$ is an eigenvector with eigenvalue $x$. Every vector is
  $\ket{w} = \sum_x \op{P}_x\ket{w}$, a sum of such eigenvectors; if $\op{A}\ket{w} = \lambda\ket{w}$,
  applying $\op{P}_x$ gives the equation $x\op{P}_x\ket{w} = \lambda\op{P}_x\ket{w}$, and
  $\ket{w}$ lies in the image of $\op{P}_\lambda$. $\op{A}$ is Hermitian because it is a linear
  combination of Hermitian operators with real coefficients. For~2: by Theorem~\ref{c03.thm.spektral},
  $\op{A} = \sum_x a_x\ketbra{\phi_x}{\phi_x}$; the projector onto the eigenspace for $\lambda$ is
  $\op{P}_\lambda = \sum_{x: a_x = \lambda} \ketbra{\phi_x}{\phi_x} = f_\lambda(\op{A})$ with
  $f_\lambda(a) = 1$ for $a = \lambda$ and $0$ otherwise, so by
  Proposition~\ref{c03.prop.funktion} it does not depend on the eigenbasis. The $\op{P}_\lambda$ are
  projectors with sum $\Id$, and $\sum_\lambda \lambda\op{P}_\lambda = \op{A}$.
\end{proof}

% vl06.tex:124 ("one therefore often reads as a definition ...").
One therefore often reads as a definition that the observables are exactly the Hermitian
operators. That is the projection-valued special case: measurements like the one of
Example~\ref{c04.ex.wirkungsgrad} are not contained in it.

\section{Time evolution}\label{c04.s2}

% Voice: vl06.tex:126-128; v08 [20:00]-[30:00].
A state characterises not a single system with somehow changing properties, but a preparation
procedure that can be set going at any time. We consider here a closed system whose conditions do
not change with time. Then it does not matter whether preparation and measurement are both carried
out somewhat later; the only thing that counts is the \emph{relative} time delay. That is an
assumption about the system, not a consequence of reproducibility: a system driven from outside in
a time-dependent way does not satisfy it.

With P0 to P4 we would in principle be finished; but because shifts in time keep turning up, we
want a convenient auxiliary postulate. It describes the reversible dynamics of a closed system
with a time-independent Hamiltonian and uses the operator $\op{U}_t$ of
Equation~\eqref{c03.eq.zeitentwicklung}.
% vl06.tex:126-130 (postulate 5, wording of v0.2) and 136 (Hermiticity added afterwards); hazard H4.3.
\begin{equation}\label{c04.eq.p5}\claim{c04.cl.p5}{author-stipulated}
  \begin{aligned}
    &\text{P5.\ \tpcifenglish{Every kind of system is further characterised by a Hamiltonian}{Jede Systemsorte ist weiter charakterisiert durch einen Hamiltonoperator}}\ \op{\Ham} , \\
    &\text{\tpcifenglish{which determines the time-evolution operator}{der den Zeitentwicklungsoperator}}\ \op{U}_t = e^{-\frac{it}{\hbar}\op{\Ham}}\text{\tpcifenglish{.}{ bestimmt.}} \\
    &\text{\tpcifenglish{If, in an experiment with a given preparation and measurement procedure,}{Wendet man in einem Experiment mit gegebenem Präparier- und Messverfahren}} \\
    &\text{\tpcifenglish{the measurement is applied with a delay of}{die Messung mit einer Verzögerung von}}\ t\ \text{\tpcifenglish{seconds, one obtains the expectation value}{Sekunden an, so erhält man den Erwartungswert}}
  \end{aligned}
\end{equation}
\begin{equation}\label{c04.eq.dynamik}\claim{c04.cl.dynamik}{derived}
  \langle \op{A} \rangle(t) = \tr\bigl(\op{U}_t\, \rho\, \op{U}_t\adj\, \op{A}\bigr) = \tr\bigl(\rho\, \op{U}_t\adj \op{A}\, \op{U}_t\bigr) .
\end{equation}
This equation follows simply from the cyclicity of the trace. It is included here because the
two sides correspond to different, but evidently equivalent, conceptions of time evolution. On the left it is natural to introduce
$\rho_t = \op{U}_t\, \rho\, \op{U}_t\adj$ as a time-dependent state and to regard the
measurements as time-independent; that is the \emph{Schrödinger picture}. For pure states this
means $\ket{\psi_t} = \op{U}_t\ket{\psi}$. On the right it is natural to introduce time-dependent
effect operators and observables $\op{A}_t = \op{U}_t\adj \op{A}\, \op{U}_t$ with a fixed state;
that is the \emph{Heisenberg picture}. Both pictures give the same probabilities; which one uses is
a matter of convenience.

One prerequisite is missing from the postulate, and we supply it here: the Hamiltonian $\op{\Ham}$
is always \emph{Hermitian}\zoomref{c04.v2}. Only then is $\op{U}_t$ unitary for all $t$, and only
then are $\rho_t$ and $\op{U}_t\adj\op{\EffOp}\,\op{U}_t$ again a state and a detector:

\begin{proposition}[Time evolution preserves states and detectors]\label{c04.prop.dynamik}\claim[proof-included]{c04.cl.dynamik-zulaessig}{derived}
  Let $\op{U}$ be unitary. If $\rho$ is a density operator, so is $\op{U}\rho\,\op{U}\adj$; if
  $0 \leq \op{\EffOp} \leq \Id$, then also $0 \leq \op{U}\adj\op{\EffOp}\op{U} \leq \Id$. The map
  $\rho \mapsto \op{U}\rho\,\op{U}\adj$ maps mixtures to mixtures and pure states to pure states.
\end{proposition}
\begin{proof}
  For every vector,
  $\melem{\psi}{\op{U}\rho\,\op{U}\adj}{\psi} = \melem{\op{U}\adj\psi}{\rho}{\op{U}\adj\psi} \geq 0$,
  and $\tr(\op{U}\rho\,\op{U}\adj) = \tr(\rho\,\op{U}\adj\op{U}) = \tr(\rho) = 1$. Likewise
  $\melem{\psi}{\op{U}\adj\op{\EffOp}\op{U}}{\psi} = \melem{\op{U}\psi}{\op{\EffOp}}{\op{U}\psi} \geq 0$,
  and the same for $\Id - \op{U}\adj\op{\EffOp}\op{U} = \op{U}\adj(\Id - \op{\EffOp})\op{U}$. The
  map is linear, and $\op{U}\ketbra{\psi}{\psi}\op{U}\adj = \ketbra{\op{U}\psi}{\op{U}\psi}$
  with $\norm{\op{U}\psi} = \norm{\psi}$.
\end{proof}

\begin{vertiefung}[id=c04.v2,
    titel={Unbounded Hamiltonians},
    wozu={In the core $\op{\Ham}$ is a Hermitian matrix, and $e^{-it\op{\Ham}/\hbar}$ is a sum over
      the eigenvectors. The most important Hamiltonians, however, are unbounded. Here is why
      ``Hermitian'' is then no longer enough, and what takes its place.},
    voraussetzung={Section~\ref{c04.s2}; blocks~\ref{c03.v5} and~\ref{c03.v7} (symmetric and
      self-adjoint, spectral theorem). The core does not need this block; Chapters~5 and~9 take
      the subject up again.},
    zurueck={For an unbounded Hamiltonian it has to say: self-adjoint. Exactly the self-adjoint
      operators generate unitary time evolutions.}]
  % Hazard H4.3: vl06.tex:138 "Only then [Hermitian] is U_t unitary"; for unbounded H symmetric is
  % not enough, self-adjoint is needed (Stone's theorem).
  For an unbounded operator, block~\ref{c03.v5} distinguishes symmetric and self-adjoint, and only
  for self-adjoint ones does the general spectral theorem,
  Theorem~\ref{c03.thm.spektral-allgemein}, hold. With it one defines
  $e^{-it\op{\Ham}/\hbar} = \op{V}\adj\, e^{-ita/\hbar}\, \op{V}$: first use the unitary operator
  $\op{V}$ of the theorem to map into $\Ltwo{M}$, then multiply by the function
  $e^{-ita(m)/\hbar}$ of modulus $1$, and map back with $\op{V}\adj$. That is a unitary operator, and the power series is
  not needed; for unbounded operators it makes sense only under very restrictive conditions. That
  it is exactly the self-adjoint operators is what Stone's theorem says:

  \begin{theorem}[Stone]\label{c04.thm.stone}\claim[proof-omitted]{c04.cl.stone}{source-quoted}
    If $\op{A}$ is self-adjoint, then $\op{U}(t) = e^{it\op{A}}$ is a strongly continuous
    one-parameter unitary group: $\op{U}(t)$ is unitary, $\op{U}(t+s) = \op{U}(t)\op{U}(s)$, and
    $\op{U}(t)\ket{\psi} \to \op{U}(t_0)\ket{\psi}$ for $t \to t_0$ and every vector. Conversely,
    every such group is of this form, with a self-adjoint $\op{A}$.
  \end{theorem}

  Reed and Simon prove the two directions as Theorems VIII.7 and VIII.8; that $\op{A}$ is fixed by
  the group follows from parts~(c) and~(d) of VIII.7, by which $i\op{A}\ket{\psi}$ is the
  derivative of $\op{U}(t)\ket{\psi}$ at $t = 0$ and $\op{A}$ is defined exactly where this
  derivative exists
  \srcref{reed-simon-1}{p.~265 and p.~266}. With $\op{A} = -\op{\Ham}/\hbar$, $\op{U}(t) = \op{U}_t$.
  A symmetric operator that is not self-adjoint generates no such group; that is the reason why
  it has to say ``self-adjoint'' instead of ``Hermitian'' for unbounded Hamiltonians.
\end{vertiefung}

% Voice: vl06.tex:136 ("An experiment with a delayed measurement is then simply an experiment with
% a different measurement ... additional piece of information about the kind of system");
% BRIEF.md:55 ("dynamics are symmetries").
An experiment with a delayed measurement is then simply an experiment with a different
measurement; in that sense P5 is compatible with the other postulates. Even so, it cannot be
proved from them, and the essential point remains: $\op{\Ham}$ is an \emph{additional} piece of
information about the kind of system. The proposition shows in what sense
dynamics is a symmetry: $\op{U}_t$ maps the state space reversibly onto itself and preserves
mixtures and pure states. Chapters~8 and~9 turn this into the general notion.

% vl06.tex:140-160 (the example is in C3 as c03.ex.qubit-zeit; here the measurement for it) and
% 162-164 (figure).  With an energy epsilon instead of H = sigma_x (migration record on vl06).
\begin{example}[A qubit under the Z measurement]\label{c04.ex.qubit-dynamik}\claim{c04.cl.qubit-dynamik}{derived}
  Let $\op{\Ham} = \varepsilon\,\pauli{x}$ with an energy $\varepsilon > 0$; by
  Example~\ref{c03.ex.qubit-zeit},
  $\op{U}_t = \cos(\varepsilon t/\hbar)\Id - i\sin(\varepsilon t/\hbar)\pauli{x}$. In the state
  $\rho = \ketbra{0}{0}$ the Z measurement after a time $t$ gives
  \begin{equation}\label{c04.eq.rabi}
    \tr\bigl(\rho_t\,\ketbra{0}{0}\bigr) = \abs{\melem{0}{\op{U}_t}{0}}^2 = \cos^2\Bigl(\frac{\varepsilon t}{\hbar}\Bigr) .
  \end{equation}
  In general a state~\eqref{c04.eq.qubit} with $\real x = 0$ stays in this plane and runs there on
  a circle about the centre $\Id/2$: with $p - \tfrac12 = r\cos\varphi$ and
  $\imag x = r\sin\varphi$,
  \begin{equation}\label{c04.eq.kreis}
    p_t - \tfrac12 = r\cos\Bigl(\varphi + \frac{2\varepsilon t}{\hbar}\Bigr) ,
    \qquad
    \imag x_t = r\sin\Bigl(\varphi + \frac{2\varepsilon t}{\hbar}\Bigr) ,
  \end{equation}
  and the Z measurement reads off the horizontal coordinate $p_t$.
\end{example}

\begin{tpcfigure}{c04.fig.messung}
  \caption{What the postulates P2, P4 and P5 fix together, for a qubit with
    $\op{\Ham} = \varepsilon\,\pauli{x}$ and the Z measurement $\op{\EffOp} = \ketbra{0}{0}$. (a) The
    states~\eqref{c04.eq.qubit} with $\real x = 0$: positivity, Equation~\eqref{c04.eq.det}, cuts
    out a disc of radius $\tfrac12$ about $\Id/2$; outside it (hatched) lie Hermitian matrices with
    trace one that are not states. The boundary consists of the pure states. $\op{U}_t$ rotates the
    disc about its centre, Equation~\eqref{c04.eq.kreis}; the circles are the orbits of the pure
    state $\ketbra{0}{0}$ (solid) and of a mixed state $\rho$ with $r = \tfrac{3}{10}$ (dotted).
    (b) The detection probability $\tr(\rho_t\op{\EffOp}) = p_t$ against time, in units of
    $\hbar/\varepsilon$: the distance $r$ from the centre in (a) is the amplitude in (b). The pure
    state oscillates between $0$ and $1$, Equation~\eqref{c04.eq.rabi}; the maximally mixed state
    $\Id/2$ always answers with $\tfrac12$ (horizontal line).}
\end{tpcfigure}

% Voice: vl06.tex:166 (closing paragraph); v08 [40:00]-[44:00].
A significant gap remains in these postulates: nowhere is it said which operators $\rho$ or
$\op{\EffOp}$ one should assign to a given experiment. That cannot be settled by a few general
postulates. The most common approach is to single out some standard observables, as a rule by
their behaviour under the symmetries of space and time. They are closely related to classical
quantities and are also named after them: position, momentum, energy, angular momentum. Whoever
wants a more precise assignment of measuring devices and operators has to analyse the measurement
process itself, especially the interaction of the system with parts of the apparatus. P0 to P5 do
not say what happens to the system after a measurement either; repeated measurements on the same
system come later. For almost all questions of this course we get by with the standard
observables; Chapter~5 begins with position.
